\chapter{A conexão entre marca sustentável e negócios na Floresta Amazônica e a percepção do empreendedor local}

\begin{flushright}
Luiz Diego Vidal Santos  \footnote[1]{Programa de Pós-graduação em Ciências da Propriedade Intelectual-PPGPI/Universidade Federal de Sergipe-UFS. *E-mail:  \href{mailto:vidal.center@academico.ufs.br}{vidal.center@academico.ufs.br}}
\\[0.2cm] 
Francisco Sandro Rodrigues Holanda \footnote[2]{Departamento de Engenharia Agronômica/Universidade Federal de Sergipe, São Cristóvão, Brasil. E-mail \href{mailto:fholanda@academico.ufs.br}{fholanda@academico.ufs.br}}
\\[0.2cm] 
Catuxe Varjão de Santana Oliveira \footnote[3]{Departamento de Ciências da Computação/Instituto Federal de Sergipe, Lagarto, Brasil. E-mail \href{mailto:catuxe.oliveira@ifs.edu.br}{catuxe.oliveira@ifs.edu.br}}
\\[0.2cm] 
Pedro Vinícius Bertulino de Menezes \footnote[2]{Departamento de Engenharia Agronômica/Universidade Federal de Sergipe, São Cristóvão, Brasil. E-mail \href{mailto:pedrovinicius18@academico.ufs.br}{pedrovinicius18@academico.ufs.br}}
\\[0.2cm] 
Alceu Pedrotti \footnote[2]{Departamento de Engenharia Agronômica/Universidade Federal de Sergipe, São Cristóvão, Brasil. E-mail \href{mailto:alceupedrotti@gmail.com}{alceupedrotti@gmail.com}}
\end{flushright}


% resumo em português
\begin{resumoumacoluna}
Os discursos sobre o empreendedorismo na região, amazônica destacam a importância de ações sustentáveis desenvolvidas por marcas locais, porém pouca atenção deu-se à visão do empreendedor quanto aos resultados das práticas aplicadas. Este estudo visa compreender a percepção dos empreendedores a partir do seu modelo de negócio relacionado à conexão da marca Amazônia. O estudo consistiu em entrevistas semiestruturadas com 8 marcas de produtos amazônicos, utilizando a teoria fundamentada. Os resultados demonstram que os empreendedores compreendem que o alcance da dimensão social depende da articulação de três fatores: dos empreendedores, dos povos originários e do conhecimento tradicional compartilhado. Observa-se que a dimensão política ainda é pouco trabalhada na região sendo pouco aplicados os seus efeitos nos produtos. Quanto à dimensão econômica, ainda é pouco favorecida pelo uso da marca sustentabilidade ambiental. Por fim, a dimensão territorial depende da criação de áreas de proteção permanente, assim como sistemas agroflorestais.

\vspace{\onelineskip}
 
\noindent
\textbf{Palavras-chave}: Floresta Amazônica; branding; Produtos sustentáveis; Intenção de compra.
\noindent

Artigo sob avaliação no periódico: Sustainable Earth Reviews %\href{https://doi.org/10.53660/CLM-106-127B}{\color{blue}{\underline{https://doi.org/10.53660/CLM-106-127B}}}



\begin{figure}[H]
  \centering
  \caption*{Resumo gráfico 4º capítulo}
  \includegraphics[scale=0.15]{Imagens/resumo_4.pdf}
  \fonte{Elaborado pelos autores, 2022}
  \label{figres:02}
\end{figure}

\href{https://empreendasustentavel.com/resumo_4.html}{\color{blue}{\underline{Download: Baixe a imagem em tamanho real}}}


\end{resumoumacoluna}
 
\section{Introdução}

A preservação da floresta amazônica é crucial para o planeta, pois representa mais da metade das florestas tropicais do mundo e um quarto de sua fauna. Especialmente para o Brasil, a floresta Amazônica se mostra ainda mais importante, à medida que a cobertura florestal desempenha um papel vital para o ciclo hidrológico, que abastece os aquíferos das regiões locais, assim como as regiões mais industrializadas da América do Sul \cite{Weng2017UnderstandingCommunities,HenrichsdeAzevedo2021FlowSystem}. Apesar da queda do desmatamento na região entre 2004 e 2012, como resultado do controle federal \cite{SilvaMontes2022TheApproach}, houve um aumento sem precedentes em 2015 e 2016. Simultaneamente, o Brasil vive a maior crise econômica de sua história recente e uma crise política, carecendo de pesquisas que possam levantar soluções plausíveis para mitigar esses impactos já percebidos \cite{Pereira2019PolicyRainforest}. Uma das formas de reduzir tais impactos é o empreendedorismo ambientalmente sustentável na localidade.

O empreendedorismo significa fazer algo novo, diferente, mudar a situação atual e buscar, incessantemente, novas possibilidades de negociações, tendo como foco a inovação e a criação de valor \cite{SocorroMarquez2021TheModel}. Outrossim, \citeonline{Leite2015} retratam o empreendedorismo como um conjunto de conceitos, métodos, instrumentos e práticas relacionadas com a criação, implantação e gestão de novas empresas ou organizações. Observando tais conceitos, pode-se entender uma empresa empreendedora como um ambiente em constante busca por soluções às fraquezas e ameaças ao ambiente, observando vantagens e oportunidades nas constantes mudanças do desenvolvimento socioeconômico \cite{Shan2018AssessingDevelopment}, gerando valor a partir da adoção de novas tecnologias e processos constatados mediante às necessidades observadas \cite{Trifonova2018EducationalRussia,Ferreira2020a,Peng2022DoesEntrepreneurship}. A relação dos empreendedores com as comunidades locais é particularmente relevante para a região amazônica, a maior floresta tropical contígua remanescente, que abriga mais de 50 milhões de pessoas, cerca de 180 comunidades indígenas com uma população de aproximadamente 440.000 povos indígenas, 1.000 comunidades quilombolas e outros povos tradicionais tais como: piaçabeiros, pescadores, catadores de açaí entre outros \cite{CIMI2017PovosEsperanca}.

As propriedades florestais globais sofrem com impacto das mudanças de hábitos da população \cite{Singh2021ForestIntervention}. Dada a extensão espacial das florestas desprotegidas, é imperativo encontrar abordagens de manejo que sustentem seus múltiplos valores ambientais, econômicos e culturais e impeçam sua conversão para outros usos \cite{Lindenmayer2012AIssues}. As empresas empreendedoras devem pensar em sustentabilidade como ato empresarial, revisitando a discussão sobre o estigma de empresas serem, por muitos, consideradas as principais geradoras de problemas ambientais e sociais, portanto, fonte da falta de atuação sustentável na sociedade \cite{Schaltegger2022ManagingField}. No entanto, enquanto essa visão tende a superestimar a possibilidade de as empresas desenvolverem os seus negócios apenas de forma exploratória, este pressuposto subestima e distorce o papel das empresas no desenvolvimento dos seus negócios sob um ecossistema formado por organizações objetivando o apoio ao empreendedorismo social e ambiental \cite{Gu2022TheCountries}, impossibilitando o empreendedorismo de ser uma possível ferramenta de promoção da sustentabilidade \cite{Breitbarth2018TheSpace}.

O empreendedorismo pode surgir também como agente mediador na promoção de desenvolvimento rural sustentável \cite{Luswaga2021Non-timberTanzania}. Ao se discutir agendas de políticas públicas promotoras de novos empreendimentos, levantando uma oportunidade de abordar a conservação ambiental e os fluxos de produção agrícola antes de serem postos em prática as atividades planejadas. Neste sentido, o empreendedorismo se torna um agente de ressignificação de métodos outrora tradicionais, convertendo-os em produtos e serviços ambientais e sociais menos impactantes ambientalmente \cite{Schaltegger2022ManagingField}, nutrindo novas formas de promoção ou pensamento entre atividades sociais e recursos ambientais \cite{Kanda2021FromEcosystems}. É importante compreender que, para entender o conceito de sustentabilidade, é necessário ter um olhar sistêmico que permeie todo o processo, envolvendo diversas dimensões, dentre as quais se destacam a econômica, a sociocultural, a político-institucional e hábitos de consumo. Neste sentido, se faz necessário compreender como o empreendedorismo adota e implementa estratégias para atingir a sustentabilidade, ou seja, ter um empreendedorismo sustentável \cite{Jarrar2022EntrepreneurshipSESOs}. O empreendedorismo sustentável ocorre no ato de desenvolver novos negócios focados em atividades de transição de consumo, relacionados com a preocupação, com o bem-estar e reconhecendo a necessidade de operar nas limitações ecológicas do planeta \cite{Munoz2018SustainableAhead}.

A contribuição do empreendedorismo para a criação de hábitos conscientes de consumo é uma questão amplamente discutida \cite{Ramadani2016AnContext}, que está atrelada ao comportamento social dos consumidores e à produção industrial. Compreender as ações que promovam um consumo mais consciente contribui para o surgimento de planos gestores empresariais mais sustentáveis, ao introduzir inovações antipredatórias que satisfazem as demandas específicas dos consumidores, enquanto lidam com a autossuficiência conscientemente duradoura sem esgotar os recursos naturais.

Um impulso para mudar os padrões de consumo dos recursos ambientais depende de modificações substanciais em toda a produção e promoção comercial dos produtos \cite{Shen2020}. Estratégias de marketing são mediadoras desse objetivo \cite{Jin2018TheDevelopment}, uma vez que o desempenho da estratégia de marketing considera impactos sociais disruptivos, uma variável dependente no feedback gerencial \cite{Bagozzi2016StateUsage}. Visando a intervenção social para conscientizar quanto à necessidade do consumo sustentável, as marcas podem se tornar ativistas de causas ao mobilizar debates em torno de uma causa política, promovendo esses debates que legitimam práticas empresariais emergentes \cite{Koch2020BrandsCase}.

Os consumidores tendem a comprar marcas  ou produtos cujas imagens coincidem com seus princípios éticos e morais, o que é um fator considerável na formação de uma relação emocional e sentimental entre um consumidor e uma marca \cite{Liu2020b}. Neste cenário, a marca é considerada o principal elemento para a compreensão da percepção do consumidor, incluindo rótulos, símbolos, descrições e/ou rurais, além de associações, resultando em um contrato flexível com o cliente para a entrega do produto ou serviço \cite{Morgan2012ARevisited,Sanchez2016RegimenPublicidad}. Tais atributos são projetados para criar consciência de identidade e construir a reputação de um produto, serviço, pessoa, lugar ou organização \cite{Sammut-Bonnici2015BrandBranding}. Como resultado, a marca se torna um importante ativo tangível e intangível para seu portador, determinando o quanto o consumidor aceita o produto da empresa signatária da marca, impactando significativamente o seu desempenho comercial \cite{Singh2021ForestIntervention}. Para que a marca se torne um importante ativo na conscientização do consumo sustentável, a gestão do negócio que busca um posicionamento estratégico com foco na sustentabilidade deve adotar um posicionamento bem definido. O \textit{Brand Equity} pode ser alcançado mediante um plano de \textit{Branding} focado na sustentabilidade, que permite à empresa manter-se competitiva e ambientalmente responsável, não perdendo lucro com os seus produtos ou serviços e seu desempenho comercial \cite{Winzar2018BrandChain}.

O papel dos empreendedores é frequentemente enfatizado nas discussões sobre desenvolvimento sustentável \cite{Larsson2016TheSystem}. Um exemplo é a iniciativa das Nações Unidas (ONU) \textit{“Apoiando Empreendedores para o Desenvolvimento Sustentável”} \cite{SEED2015TurningEntrepreneurship}. Porém, para os empreendedores que já contam com uma atividade econômica iniciada, a transição do estágio de atividade de exploração para um estágio sustentável, se mostra difícil \cite{Imran2019ResourceCommitment}. Nesse estágio empresarial, a maioria das empresas já consolidadas negligencia os procedimentos após atingir o estágio de sustentabilidade que considera um desempenho percebido em lucros, vendas e satisfação do cliente, mas não em sustentabilidade \cite{Ismail2022EntrepreneursInnovations}. Em outros termos, para empresas menores, esta transição é um grande desafio enfrentado quando se trata de construção de uma marca pautada na sustentabilidade. Diante disso, a maioria das empresas, especialmente as que operam em regiões pouco populosas, declinam ou liquidam estes planos gestores quando atinge a maturidade e entram no mercado competitivo restrito e pouco preocupado com a sustentabilidade ambiental \cite{GalanMashenene2020PerformanceView}. 

Tradicionalmente, estratégias para sustentabilidade ambiental nas empresas é medida pelo lado da demanda. Controle de recursos materiais e medidas de progresso financeiro têm sido aplicados por meio de métricas relacionadas aos negócios, como lucratividade, geração de renda e receita. Todavia, uma das formas de compreender as atitudes assertivas da transição para a sustentabilidade ambiental é compreender os traços de personalidades empresariais que mediam a relação entre as motivações empreendedoras e a relação com a sustentabilidade ambiental \cite{Ismail2022EntrepreneursInnovations}.


Com base nisso, neste artigo se objetiva compreender a percepção dos empreendedores a partir do seu modelo de negócio referente à conexão da marca à Amazônia e seus aspectos de exploração sustentável de seus recursos naturais. 

\section{Material e métodos}

\subsection{Contexto de pesquisa}

O presente estudo estabelece como pressuposto um modelo empiricamente fundamentado que visa compreender como os empreendedores interpretam as dimensões do desenvolvimento sustentável na região amazônica, em especial as questões norteadoras das dimensões para o desenvolvimento sustentável. Para atingir este objetivo, foi utilizada a teoria fundamentada em dados \cite{Strauss1987QualitativeScientists,Strauss1990BasicsResearch}. A teoria fundamentada em dados (\textit{Grounded Theory}) é particularmente aplicável em estudos exploratórios qualitativos, sendo frequentemente usada para estudar comportamentos sociotécnicos em domínios emergentes de pesquisa \cite{Foley2015UsingExperiences}. Tal método descreve as seguintes fases para a sua condução: (i) organização da análise; (ii) codificação aberta; (iii) codificação axial; (iv) codificação seletiva e (v) delimitação da teoria.

Como fonte de dados, realizaram-se entrevistas com proprietários de 8 (oito) marcas de produtos amazônicos, que desenvolvem seus produtos considerando a sustentabilidade florestal. O estudo concentrou-se em 6 cadeias socio-produtivas da região: Turismo, Gastronomia, Biotecnologia, Sistemas Agroflorestais, Produtos Florestais Não Madeireiros e Moda Sustentável.

%\begin{figure}[H]
 % \centering
 % \caption{Mapa geográfico da região da Amazônia no Brasil e a localização das marcas dos empreendedores participantes}
%  \includegraphics[scale=1]{Imagens/fig_6.pdf}
%  \fonte{Adaptado de \citeonline{IBGE2020}}
%  \label{fig:6}
%\end{figure}

\subsection{Fase de organização}
\subsubsection*{Desenvolvimento de roteiro de entrevista}

Na primeira fase, foi elaborado um roteiro de entrevista semiestruturado (Apêndice B). As entrevistas com os proprietários das marcas foram realizadas por videochamada no período de 15 defevereiro 2023 a 16 junho 2023. O instrumento inicial tem como base as dimensões da sustentabilidade propostas por \citeonline{Sachs2002CaminhosSustentavel}. Para as dimensões: econômica (fonte de matéria-prima, produtos, etc.), social (paridade, masculino feminino, conscientização de comunidades locais, etc.) e ambiental (redução de embalagens, redução das emissões de $CO_2$ etc.), foi adaptado o questionário proposto por \citeonline{Postel2006LaFordiste} \textit{Triple Bottom Line}. As dimensões: cultural, ecológica, territorial, política (nacional) e política (internacional) foram avaliadas utilizando adaptações do questionário proposto por \citeonline{Grigorescu2019KeySustainability}.

De modo a atender os princípios éticos em pesquisa com seres humanos, o estudo foi encaminhado ao Comitê de Ética em Pesquisa com código CAAE: 55999722.1.0000.5546, conforme a Resolução do Conselho Nacional de Saúde 466/2012. Os participantes que aceitaram participar do estudo foram esclarecidos quanto aos objetivos, a coleta de dados, bem como, as informações contidas no Termo de Consentimento Livre e Esclarecido (TCLE). Mantendo a segurança ética da pesquisa, após a obtenção do consentimento dos participantes, as entrevistas foram desidentificadas, gravadas e em seguida, transcritas.

A adaptação transcultural do instrumento inicial foi composta por três fases: a primeira foi realizada logo na tradução do instrumento, sendo realizada por dois tradutores independentes. No segundo momento, foi realizada a validação dos itens por meio do Coeficiente de Validade de Conteúdo (CVC), técnica que envolve a busca do bom senso de uma comunidade de juízes (ou seja, profissionais efetivamente envolvidos na área de pesquisa) \cite{Gurgel2021}.

Na terceira fase, foi realizada um \textit{back-translation}, validada pelos tradutores, buscando verificar se a nova proposta do questionário condizia com os objetivos propostos nos questionários originais. O roteiro de entrevista final foi consolidado em 22 (vinte e duas) questões que abordaram as 8 dimensões do desenvolvimento sustentável proposto por \cite{Sachs2002CaminhosSustentavel}. As entrevistas duraram em média de 30 a 45 minutos. 

\subsection{Grounded Theory}

Para a análise dos dados, uma abordagem de \textit{Grounded Theory} foi adaptada para corresponder ao objetivo do estudo, que inclui quatro etapas: codificação aberta, codificação axial, codificação seletiva e delimitação da teoria \cite{Glaser2017DiscoveryResearch, Mills2006TheTheory}.

\subsubsection*{Codificação aberta}

Durante a fase de codificação aberta (i), deu-se início ao processo de comparação das ocorrências textuais dos questionários. Análises dos dados foram realizadas, por meio do auxílio do software de R com pacote IRaMuTeQ. Nesta pesquisa cada resposta obtida foi considerada uma Unidade de Contexto Inicial (UCIs).

\subsubsection*{Codificação axial}

A fase de codificação axial (ii) foi realizada por Classificação Hierárquica Descendente (CHD), onde os segmentos de texto são classificados em função dos seus respectivos vocabulários, resultando em classes lexicais. Foram considerados termos, palavras classificadas como substantivos, adjetivos e verbos. Demais conjuntos e termos presentes nos Seguimentos de Textos (STs) foram excluídos da CHD. 

As classes lexicais obtidas nesta fase foram utilizadas como códigos construtores do guia. Por fim, foi realizada uma análise de similitude que identifica as ocorrências entre as palavras, trazendo indicações da conexidade entre as palavras. Ao final, foi possível identificar quais os principais códigos que norteiam a percepção dos empreendedores quanto à sustentabilidade ambiental da região da Amazônia.

\subsubsection*{Codificação seletiva}

A fase de codificação seletiva (iii) focou no desenvolvimento do guia de códigos. Para o desenvolvimento dos \textit{clusters} presentes no guia de códigos, foram utilizadas inicialmente as classes emergidas do CHD. As classes foram analisadas individualmente por três pesquisadores antes de considerar todas as transcrições como um conjunto de dados completo. 

As entrevistas foram codificadas e memorizadas por cada membro da equipe de pesquisa. Partimos do pressuposto de que o trabalho continha inicialmente cinco códigos relativos às classes emergidas da CHD. 


\subsubsection{Delimitação teórica}

A fase de delimitação da teoria (iv) foi desenvolvida para o guia de códigos e a associação das respostas com os códigos observados. Esta fase foi desenvolvida para dois estudos que não foram publicados, mas foram codificados usando o software Atlas.ti 8 \cite{Paulus2019ItMethod}.

\section{Resultados e discussão}

Nesta sessão serão fornecidas evidências a partir das percepções de empreendedores locais da Amazônia sobre a conexão da sua marca com a sustentabilidade ambiental florestal, seu papel como empreendedor e suas metas de gestão para alcance da sustentabilidade.

O corpus textual estudado foi constituído pelas respostas dos entrevistados contendo 12 textos (questões), separados em 601 segmentos de texto (ST), com aproveitamento de 481 STs (80,55\%). Emergiram 20.994 ocorrências (palavras, formas ou vocábulos), sendo 2.926 palavras distintas e 1.536 com uma única ocorrência. A partir dos agrupamentos dos seguimentos de textos, foram consideradas 5 classes a partir do agrupamento e classificação hierárquica descendente CHD. 

O conteúdo analisado foi categorizado em cinco classes: Classe 1 - \textit{“Agentes e comunidades envolvidas” - (ACI)}, com 101 ST (20.73\%); Classe 2 - “Comercialização dos produtos”, com 106 ST (22.20\%); Classe 3 - “Disponibilidade de recursos e matéria-prima utilizada”, com 77 ST (16.34\%); Classe 4, “\textit{Brand equity} percebido” relacionado à floresta, com 115 ST (23.9\%); e Classe 5 - “Visão das atividades florestais sustentáveis”, com 81 ST (16,83\%).

\subsection{Resultados de codificação axial e codificação seletiva}

A Classe 1 - \textbf{“Agentes e comunidades envolvidas (ACI)”} compreende 19,02\% (\textit{f} = 209 ST) do corpus total estudado, foi composto por palavras e radicais no intervalo entre $\chi_2$ = 2,11 (aspecto) e $\chi_2$ = 38,33 (coisa). Nesta classe, surge conteúdos referentes à percepção dos empreendedores sobre a participação dos povos originários e o envolvimento na sustentabilidade ambiental da floresta. Além disso, sinaliza aspectos do desenvolvimento de produtos negociados pelas marcas estudadas e a responsabilidade empresarial destas que tange à manutenção sustentável da floresta circularmente, direcionado principalmente ao compartilhamento dos conteúdos relativos à sustentabilidade. Alguns segmentos de textos caracterizam essa classe:

\textit{“(...) Às vezes, como as pessoas das comunidades são ribeirinhas, indígenas, não têm noção do valor econômico e cultural dos produtos que têm à disposição.”} (empreendedor 01).

\textit{“(...) O nativo, é empreendedor, é microempreendedor, mas precisa ter esse empoderamento, senão as coisas não dão certo. Por exemplo, a sociedade pode ver aquele empreendedor como semianalfabeto, mas ele é um agente de alta inovação.”} (empreendedor 04).

\textit{“(...) Estamos tão preocupados em melhorar a qualidade de vida da comunidade, das mulheres envolvidas que trabalham para a comunidade, o impacto ambiental que temos lá, socioambiental, mas também com os produtos que usamos, que também impacta o meio ambiente.”} (Empreendedor 03).

Classe 2 - \textbf{“Comercialização de produtos e mercado consumidor” — (MPCM)}, compreende 2.02\% (\textit{f} = 9 ST) do corpus total. Consiste em palavras e radicais na faixa entre $\chi_2$ = 2.73 (Brasil) e $\chi_2$ = 37.48 (preço). Consiste em segmentos de texto relacionados à percepção de seus produtos e preço justo percebido pelos empreendedors. Alguns textos podem explicar tais segmentos:

\textit{“(...) O preço justo deve atender a todos: cliente, comunidade, minha empresa e meio ambiente.”} (empreendedor 02).

\textit{“(...) Até eu me surpreendo e entendo por que é um produto mais caro que produtos similares sem ser sustentável, mas o resultado é maravilhoso. A qualidade do produto e a responsabilidade ambiental, a meu ver, é fundamental para poder vender a esse preço”}. (empreendedor 05).

É importante destacar que, para os empreendedors participantes da pesquisa, a redução dos riscos da aceitação de preços justos está atrelada a um esforço constante de se apresentar de forma sustentável perante a sociedade, utilizando certificações e discursos perante a comunidade, mostrando a importância de posicionamento da marca por meio de letreiros que evidenciam suas práticas ambientalmente sustentáveis \cite{Dwivedi2018BrandUniqueness}.

Classe 3 - \textbf{“Disponibilidade de recursos e matéria-prima utilizada” — (ARM)}, compreende 16,34\% (f = 67 ST) do corpus total estudado. Consiste em palavras e radicais na faixa entre $\chi_2$ = 2,07 (fornecedor) e $\chi_2$ = 43,00 (matéria-prima). A classe 3 possui segmentos relacionados ao nível de disponibilidade e a aquisição da matéria-prima utilizada em seus produtos, bem como a responsabilidade da marca em manter-se junto às comunidades extrativistas, buscando com isso um comércio justo.

\textit{“(...) Elas falam que não adianta porque não tem demanda e a gente fica tentando mostrar que existe demanda, assim se elas oferecerem, no momento de safra, assim a gente tem condições de comprar numa escala maior”}. (Empreendedor 03).

\textit{“(...) Eles nos convidaram para ir lá dar uma oficina para eles, e ensinar fazer a tinta que tenha durabilidade maior para eles usarem”}. (Empreendedor 06).

Para criar novos negócios em áreas florestais, é importante que diferentes pessoas entendam como as coisas funcionam, usando o que sabem \cite{Nayha2019TransitionSustainability}. A importância de discutir e compartilhar conhecimentos ecológicos tradicionais para a sustentabilidade fica evidente quando se percebe que tais conhecimentos, transmitidos de geração em geração, tratam essencialmente de como garantir o uso contínuo de recursos naturais limitados \cite{Kim2019TraditionalKorea}. 

Em outras palavras, entender o conhecimento tradicional, convertendo-o em tecnologias sustentáveis, fornecem novas maneiras de manter o valor dos serviços ecossistêmicos e do meio ambiente conservados, pois as definições de sustentabilidade e desenvolvimento sustentável também podem ser baseadas em múltiplas perspectivas sociais, valores morais e crenças, que influenciam como uma determinada sociedade constrói a realidade em um determinado momento \cite{Easdale2020StrainsRegion}. No entanto, é importante lembrar que, para o conhecimento tradicional permanecer disponível nas comunidades locais, é necessário que os membros da comunidade consigam transmiti-lo eficazmente e de maneira segura, exigindo certa confiança entre as partes \cite{Berkes2000RediscoveryManagement}. Essa confiança pode ser obtida através da troca de informações e conhecimentos.

A classe 4 - \textbf{“Brand equity percebido” — (BEP)} e signos relacionados à floresta amazônica”, compreende 23,09\% (\textit{f} = 98 ST) do corpus total da conclusão 4. Consiste em palavras e radicais na faixa entre $\chi$ = 2.07 (Design) $\chi_2$ = 52.86 (Usar). A classe 4, um volume maior de palavras da ST, consiste em segmentos de texto que retratam a percepção do empreendedor sobre a operacionalização, as estratégias visuais e posicionamentos voltados para a sustentabilidade ambiental.

\textit{“(...) A nossa logo é um macaco, da região.”} (Empreendedor 07).
\textit{“(...) A nossa marca é o nome de uma papa feita de farinha produzida na Amazônia, consumida pelos indígenas, tem como finalidade dar substância. Aparentemente parece uma coisa muito simples, água e farinha, mas para você ter a farinha você deve ter todo um processo e cuidado.”} (Empreendedor 06).

Para marcas que operam posicionamentos voltados para a sustentabilidade ambiental, é essencialmente importante uma gestão estratégica que apresente suas perspectivas e preocupações pelo meio ambiente através dos signos adotados em seu discurso \cite{Osterwalder2019}. Uma vez que as atividades sustentáveis influenciam positivamente no seu desempenho e fornecem percepções significativos, a compreensão dessas atividades é geralmente percebida pelo simbolismo imbuído dos signos da marca, pois, para compreender e comunicar significados complexos sobre o valor, condição e moralidade da empresa e os consumidores da mercadoria usam o uso dos meios identificáveis aplicados pela empresa \cite{Yu2021CanPremium}.

A Classe 5 compreende 16,83\% (f 64 ST) do corpus total considerado. Composto por palavras e radicais no intervalo entre $\chi_2$ = 2.47 (região) $\chi_2$ = 42.47 (forma). Nesta classe 5 \textit{“Visão das atividades florestais sustentáveis — (VSFA)”}, é demonstrada a percepção dos empreendedores quanto à preservação dos passivos ambientais relacionados às suas áreas de atuação na floresta e seus vínculos psicológicos. Os seguintes STs descrevem e explicam esta classe.

\textit{“(...) Por exemplo, eu fico muito triste aqui em Belém porque existem vários canais que eram chamados de Veneza e hoje em dia esses canais, como a Doca que fica em um bairro nobre, se tornou um esgoto ao ar livre.”} (Empreendedor 07).

\textit{“(...) Então hoje a quem estamos guardando a floresta? Vejo que  que já foi tudo desmatado e na realidade temos uma superfruta que nessa região se reproduz naturalmente e as pessoas continuam queimando e matando elas.”} (Empreendedor 05).

Os dados demonstram que as alterações ambientais existentes, tais como: assoreamentos e interrupções dos rios e igarapés, desmatamentos acentuados, contaminações do solo, geram preocupações para os empreendedores e impactam negativamente em seus negócios. 

Com base nas cinco classes elencadas e na perspectiva da teoria fundamentada, foi possível determinar um número expressivo de códigos convergentes para os grupos e com cinco classes distintas. A Figura 2 mostra as classes do Dendrograma, elaborado a partir da Classificação Hierárquica Descendente — CHD.

\begin{figure}[H]
  \centering
  \caption{Dendrograma das classes DHC}
  
\includegraphics[scale=0.4]{Imagens/fig_7.pdf}
  \fonte{Elaborado pelos autores, 2022}
  \label{fig:7}
\end{figure}

A CHD formou dois subgrupos, um grupo contendo as classes 2 e 3, abrangendo temas relacionados à comercialização, produtos e serviços desenvolvidos, e as classes 1.4 e 5, que estão mais relacionadas à percepção de novos empreendedores na região amazônica, dada a importância da ligação da sua marca com a exploração sustentável da floresta. Esta organização de subclasses vai se expressar na fase de codificação seletiva. 

Com base em subgrupos e códigos contidos no guia, é possível entender quais fatores levaram os empreendedores a atuarem de forma sustentável na floresta amazônica e em que medida esses fatores podem beneficiar ou afetar negativamente as relações de longo prazo entre marcas e a sustentabilidade ambiental da Amazônia brasileira. A Tabela \ref{tab:7} apresenta os principais códigos observados. Buscando facilitar o entendimento, a partir da fase de codificação seletiva, assumimos o termo \textit{cluster} como os conjuntos de códigos pertencentes às classes iniciais que imergiram na fase axial.

\newpage
\begin{landscape}
\begin{longtable}{@{}ccp{4cm}p{9cm}@{}}
\caption{Códigos e guias de consulta}
\label{tab:7}\\
\toprule
 &
  \textbf{Iniciais} &
  \multicolumn{1}{c}{\textbf{Código}} &
  \multicolumn{1}{c}{\textbf{Descrição}} \\* \midrule
\endfirsthead
%
\endhead
%
\bottomrule
\endfoot
%
\endlastfoot
%
\multirow{4}{*}{\textbf{\rotatebox{90}{\textit{CLUSTER} 1}}} &
  1A &
  Compreensão os agentes envolvidos &
  O código “Agentes envolvidos” neste \textit{cluster} refere-se aos agentes locais envolvidos nas atividades sustentáveis florestais, desde o extrativismo, até a aquisição dos produtos. \\
 &
  1B &
  Entendendo como as comunidades se envolvem &
  O código compreenderá as comunidades tradicionais envolvidas na marca, tais como: indígenas, quilombolas e ribeirinhos. \\
 &
  1C &
  Ação governamental &
  O código objetiva perceber a atuação estatal e de ONGs frente à sustentabilidade ambiental \\
 &
  1D &
  Percebendo o impacto social &
  O código objetiva a perceber o impacto social da marca frente às florestas. \\
\multirow{4}{*}{\textbf{\rotatebox{90}{\textit{CLUSTER} 2}}} &
  2A &
  Buscando fonte de fornecimento da Matéria-prima &
  Sobre quais as fontes de aquisição da matéria-prima utilizada e quais os agentes produtos ou extrativistas. \\
 &
  2B &
  Produzindo meios sustentáveis de manutenção florestal &
  Sobre quais as articulações as marcas desenvolvem ou buscam desenvolver para a manutenção sustentável da floresta. \\
 &
  2C &
  Reconhecendo a importância do conhecimento tradicional &
  O código neste \textit{cluster} visa obter qual a visão frente ao conhecimento tradicional das comunidades tradicionais e o impacto no empreendedorismo. \\
 &
  2D &
  Negociação de produto/serviço comercial &
  O código neste \textit{cluster} compreende quais os principais produtos/serviços estão sendo trabalhados de modo sustentável na floresta Amazônica do Brasil e estudados nesta pesquisa. \\
\multirow{4}{*}{\textbf{\rotatebox{90}{CLUSTER 3}}} &
  3A &
  Disponibilizando os recursos produzidos &
  Visa compreender em que medida os recursos utilizados pelas marcas para desenvolvimento de seus serviços/produtos, estão sendo disponibilizados, tanto para eles quanto para os demais envolvidos no ecossistema negocial local (concorrentes). \\
 &
  3B &
  Entendendo os Processos &
  Sobre quais os processos de produção e extrativismo utilizados pelos agentes produtores, e quais os impactos gerados e de que forma são minimizados. 
  \\
 &
  3C &
  Buscando dimensionamento estratégico &
  Sobre qual o perfil do mercado consumidor e qual escala visa alcançar sem impactar negativamente a floresta Amazônica. 
  \\
 &
  3D &
  Buscando o comércio justo &
  Visa compreender a margem justa de preço para negociação de seus produtos considerando as dificuldades da atuação sustentável e seus entraves para entrada no mercado 
  \\
    \multirow{3}{*}{\textbf{\rotatebox{90}{\textit{CLUSTER} 4}}} &
  4A &
  Planejando a marca &
  O código neste \textit{cluster} refere-se às estratégias de \textit{Branding} adotadas pela empresa visando o posicionamento estratégico frente ao Brand “sustentabilidade ambiental” \\
 &
  4B &
  Delineando os sinais distintivos &
  Sobre quais os signos as marcas utilizam relacionados à floresta, ou à comunidade, visando alcançar a associação dos produtos com a floresta, pelos consumidores \\
 &
  4C &
  Posicionando-se estrategicamente &
  Sobre qual o mercado a marca visa alcançar (regional, nacional ou internacional) \\
\multirow{3}{*}{\textbf{\rotatebox{90}{\textit{CLUSTER} 5}}} &
  5A &
  Buscando a sustentabilidade ambiental &
  O código neste \textit{cluster} refere-se a qual o objetivo principal a empresa tem em vista alcançar, quanto à sustentabilidade ambiental da floresta, tais como: preservação dos conhecimentos tradicionais, preservação das espécies em seu habitat ou a manutenção dos povos originários frente ao seu local, tradição e cultura. \\
 &
  5B &
  Observando as ocorrências deletérias &
  As principais atividades de degradação observadas pelas marcas em sua localidade \\
 &
  5C &
  Percebendo a sustentabilidade como um fator positivo &
  O código neste \textit{cluster} visa a compreensão da percepção frente ao futuro da floresta e as atitudes sustentáveis que são observadas e que geram resultados positivos \\* \bottomrule
\end{longtable}
\end{landscape}
\newpage

\subsection{Delimitação da teoria a partir das dimensões do desenvolvimento sustentável}

A fase de delimitação permite compreender mais sistematicamente os fatores que levaram o surgimento dos \textit{clusters} na fase axial, por meio dos \textit{clusters} de códigos baseados nas classes da codificação axial. 

\subsubsection*{Dimensão Social}

Ao analisar os Códigos 1A e 1B presentes no \textit{Cluster} 1, é possível observar um encadeamento de ideias que constroem uma visão ligada a dimensão social da sustentabilidade. Nesse contexto, a dimensão social deve ser entendida como a busca da máxima igualdade para os sujeitos envolvidos, a fim de construir uma população onde a distribuição dos bens (materiais e intelectuais) seja mais equânime, buscando, assim, melhorar significativamente os direitos e as condições \cite{Quimby2018ParticipationComanagement}. 

\textit{“(...) Sim. Sempre que precisamos de expressar algo nas embalagens a gente acaba tomando emprestado deles, das comunidades, com lógico, devido consentimento de uso. A gente até pede sugestões para eles de algum grafismo, algum tipo de arte é que eles usam para gente incorporar isso para dentro da nossa marca”}. (Empreendedor 2).


Partindo do entendimento de que todos os agentes envolvidos na atividade sustentável são de suma importância para a construção de um plano estratégico de sustentabilidade, e que tais ligações denotam envolvimento da marca com a preservação ambiental local, os agentes envolvidos também determinam o cumprimento pleno da dimensão política nacional e da dimensão social \cite{Sachs1993EstrategiasXXI}. A sustentabilidade política deve ter contribuição não só da comunidade local, mas é preciso mobilizar toda a sociedade, englobando o papel do governo, das instituições, do empresariado e abrangendo o que muitos autores chamam de sustentabilidade institucional nesta dimensão. Essa heterogeneidade e complexidade de práticas sociais que permeiam as paisagens devem ser reconhecidas para integrar diversas medidas de sustentabilidade florestal \cite{Gebara2019DeconstructingApproach}. Nesse contexto, a proximidade dos empreendedores locais aos povos das comunidades tradicionais é importante para o consumo dinâmico dos recursos naturais disponíveis na floresta. Uma vez que a dinâmica econômica predominante na região ignora a maioria a população local, causando danos ambientais e sociais, como desmatamento e degradação de recursos, colocando em risco a subsistência de inúmeras pessoas dependentes da floresta \cite{Pokorny2013FromPoverty,Godar2014Actor-specificAmazon}.

\subsubsection*{Dimensões políticas (inter)nacionais}

Quando questionados sobre as ações de agentes governamentais que determinam a dimensão política nacional (1C) sobre sua parcela de importância para desenvolvimento do empreendedorismo sustentável, os participantes, em sua maioria, afirmaram não serem assistidos, e que a política de assistencialismo somada a negligência dos órgãos públicos dificultam suas atividades locais ou desestimulam a aceitação de futuras parcerias. 

\textit{“(...) Como eu te explico? Fiquei cansada da política de assistencialismo, entendeu? O assistencialismo me cansou, queria ver além daquilo.”} (Empreendedor 1)

\textit{“(...) Nesse sentido de atuação estamos muito distantes disso. Assim, é muito distante, e a percepção das pessoas locais com relação a isso ainda não é ótima, sempre imagem que é apenas assistencialismo”}. (Empreendedor 04).

\textit{“Aqui na região não têm muitas ações, eu conheço uma ONG que agora está com um trabalho para educação com um grupo local, também tem uma empresa pública que está com um trabalho de recuperação de manguezais. Mas eles estão do outro lado da região. De órgãos públicos prefiro não falar, porque assim mais atrapalha que ajuda a gente aqui”}. (Empreendedor 05). 

Atrelada à baixa atuação governamental percebida por eles, está a preocupação com a visão dos consumidores no que tange à sustentabilidade regional aplicada pelos órgãos públicos, e o quanto isso impacta em seus negócios.

\textit{“(...) Aqui no Tapajós, a gente percebeu que no feriado de Carnaval, com certeza teve uma redução do turismo por conta desse noticiário nacional das águas do Tapajós estarem turvas e tudo por conta da mineração ilegal e tudo mais.”} (Empreendedor 03). 

\textit{“(...) É que a gente fica meio que de mãos atadas e tentando ver soluções quanto a questão do lixo deixado na Floresta nacional do rio Tapajós. (...) Como as comunidades ainda vivem com muito lixo, não existe uma coleta seletiva de lixo, para se desfazer eles queimam. Ai quando você chegar lá, você fica, meu Deus, gente, tem que trazer uma coleta seletiva, o governo tem que vir trazer alguma forma de coleta. A gente então tenta que conscientizar de alguma forma.”} (Empreendedor 01).

O rio Tapajós é originalmente caracterizado como um rio de águas cristalinas (água claras), com baixa concentração de sedimentos em suspensão \cite{Diniz2022NotaTapajos}. Em pesquisas atuais, foram observados valores altos de Hg, tanto relacionados com a mineração no rio Tapajós e seus afluentes \cite{SilvaMontes2022TheApproach} quanto relacionados à lixiviação do solo erodido após o desmatamento e a queima da cobertura florestal \cite{Lino2019TotalAmazon}.


\subsubsection*{Dimensão psicológica}

Analisando o código 1C, é possível perceber uma ligação entre ele e os códigos 2D e 3A, à medida que visualizamos uma afinidade, além do econômico, pela floresta e seus recursos. É possível notar uma ligação hereditária dos entrevistados para com a floresta e o quanto isso impacta na busca pela preservação da floresta.

\textit{“(...) Sou filho de mãe quilombola e desde muito cedo aprendi a conviver com os conhecimentos tradicionais, somando aos conhecimentos que adquiri com minha formação, vi que a agregação de valor poderia existir nos meus produtos e ainda representar e preservar a cultura e identidade dos povos que vivem na Amazônia”}.(Empreendedor 02).

A dimensão psicológica relaciona-se diretamente com o bem-estar mental. Nesse contexto, a definição de bem-estar mental está relacionada ao aumento do pertencimento pessoal, satisfação e felicidade tendo significados maleáveis, permitindo a cada um tomar conhecimento da realidade que o cerca, por meio das diferentes experiências, expectativas, motivações tanto individuais quanto para com a comunidade \cite{Mogilner2012HowChoice}. 


\textit{“(...) Na realidade, o grande feito quem fez, não foi meu Pai. Ele recebeu dos meus avós, a partir de um tempo ele deixou uma área só conservada e estou seguindo o legado dele”}.(Empreendedor 05).

\textit{“(...) Quando mudei para cá, gostava de levar as pessoas para conhecer as belezas naturais dessa região que cresci. (...) assim ponderei fazer um turismo fora do convencional, aquele de sol e praia, assim comecei um serviço diferente somado a minha paixão pedalar, tudo somou positivamente”}. (Empreendedor 07).

\textit{“(...) Temos feito esse trabalho (educação) com outras comunidades que trabalham com matérias-primas passíveis de uso nos nossos produtos, e que tem potencial de uso sem degradar. E isso tem sido ótimo para todos nós”}. (Empreendedor 06).

\citeonline{Wang2022TheCulture} propõe que o bem-estar, componente da dimensão psicológica da sustentabilidade, pode ser obtido também pela transformação os desejos individuais através da construção de uma unidade do homem com a natureza e a sociedade. Este conceito se expressa inconscientemente nos empreendedores entrevistados, à medida que visualizamos uma sensação nostálgica ao falar da relação familiar com a região amazônica e as comunidades, demostrando uma relação positiva entre a floresta, as comunidades e as histórias de vida dos empreendedores.  

\subsubsection*{Dimensão cultural}

Os códigos presentes no \textit{Cluster} 2 do guia formam uma percepção acerca da dimensão ambiental e a relação com o acesso às fontes de matérias-primas utilizadas no desenvolvimento de seus produtos, os pontos de extração e a dimensão cultural ligada à importância dos conhecimentos tradicionais necessários para a existência de uma economia circular. 

Essa percepção remete à importância da presença das informações disponíveis pelos nativos, seus saberes e costumes tradicionais, muito presentes nos povos originários da Amazônia \cite{Study2021TheChallenges}. Os códigos do \textit{Cluster} 1 (1B e 1C) emergem dessa concepção de negócio. Assim como os achados encontrados por \cite{Shafqat2022AdaptingPakistan}, o conhecimento dos povos tradicionais locais são algumas das descobertas bastante significativas para formular a percepção de importância da partilha de conhecimento para com a comunidade e os empreendedores. 

Os argumentos que os empreendedores desta pesquisa apresentaram de base para valorização do compartilhamento com os originários estão baseados nos conhecimentos ancestrais, e as soluções que eles fornecem para uma habitabilidade sustentável a partir de uma lente única forma fundamentais para formular os índices do \textit{Cluster} 2. 

\textit{“(...) É, na verdade, uma coisa que a gente precisa ir lá aprender com eles, e esse é um aspecto que a gente precisa aprender. Eles sabem como gerenciar florestas incrivelmente. Então, às vezes dá vontade de chorar de emoção.”} (empreendedor 01).

\textit{“(...) Exatamente é que as pessoas não entendem que isso é importante, mas é uma sensibilidade diante da vida. Não é uma sensibilidade de uma troca financeira. Não querem o dinheiro, não querem nenhum bem material. Eles esperam uma troca de conhecimento, uma demonstração de respeito.”} (Empreendedor 02).

\subsubsection*{Dimensão ambiental}

A sustentabilidade ambiental da Floresta Amazônica depende não apenas da aplicação de políticas tradicionais, mas também da responsabilidade ambiental e da troca de conhecimentos com os povos tradicionais  \cite{Postel2006LaFordiste}. A orientação e conscientização dos povos originários sobre a sustentabilidade ambiental também percebida pelos empreendedores como ferramentas eficazes na busca da sustentabilidade.

\textit{“Uma coisa que a gente se preocupa muito é quando a gente vai fazer as nossas visitas. É que a gente fica meio que de mãos atadas e tentando ver soluções quanto a questão do lixo deixado na Floresta nacional do rio Tapajós. Mas assim, as comunidades ainda não têm ideia de como descartar aquele lixo não degradável. Como eles tratam o lixo impacta na beleza do lugar, que influenciam, no que o turista pensa sobre aquele lugar. Eles têm ainda, infelizmente, uma relação com aquele lixo como se fosse tudo biodegradável. Só que para a gente, né, que vem de outro tipo de educação, a gente sabe que aqueles lixos ali vão perdurar por centenas de anos, então tentamos ajudá-los orientando algumas formas de coleta seletiva.”} (Empreendedor 01).

As informações quanto à dimensão da sustentabilidade ambiental podem ser planejadas por meio da execução e do monitoramento de atividades de restauração através de estratégias educativas que utilizem efetivamente os recursos naturais de forma sustentável, alocados por meio de programas de educação natural no campo \cite{Reyes-Garcia2019TheRestoration}. Em alguns estudos, os povos originários foram responsáveis pelo monitoramento ambiental, construindo relações com plantas, terras e até mesmo pela restauração ecológica \cite{Wehi2017ImportanceRestoration,Thompson2020AMonitoring}. 

\subsubsection*{Dimensão econômica}

Os códigos do \textit{Cluster} 4 estão focados na dimensão econômica, à medida que descreve a compreensão dos empreendedores sobre a sua capacidade de sx'x', expansão das vendas para o mercado internacional, por meio das estratégias do posicionamento da marca. Num ambiente de concorrência acirrada e barreiras geográficas, a introdução de novos métodos de \textit{Branding}, especialmente àqueles relacionados à biodiversidade e proteção ambiental, promove a abertura para novos mercados e a conquista de um valor na marca \cite{Gupta2016EnhancingRecycling}.

Ainda sobre a dimensão econômica, podemos perceber nos códigos 3C e 3D, a busca pelo comércio justo para todos os agentes envolvidos, considerando as dificuldades de se empreender sustentavelmente, comportamento que reforça os dados encontrados na fase axial. Os planos de comércio justo são formas proeminentes de governança nova ou privada que, se bem tratadas e gerenciadas, podem contribuir não apenas para o desenvolvimento sustentável, mas também para o esverdeamento da economia \cite{Ribeiro-Duthie2021FairEconomy}.

\textit{“(...) Até fui atrás do selo de certificação de comércio justo.”} (empreendedor 03).

\textit{“(...) Sim, mas é preciso trabalhar a imagem porque em inúmeras feiras internacionais que fomos, constatamos que o Brasil não tem mais a credibilidade de antes. É triste, é muito triste.”} (Empreendedor 02).

\textit{“(...) Temos nos questionado visto que muitas pessoas fazem ou vendem produtos dizendo ser produtos da Amazônia, mas apenas a matéria-prima vem da Amazônia, extraída de qualquer forma e a um preço barato. Agrega valor lá fora, dizendo ser daqui da Amazônia.”} (Empreendedor 05).

O Comércio Justo cria redes alternativas concretas de produtos e, simultaneamente, promove normas e práticas alternativas nas arenas comerciais, possibilitando a entrada de novos empreendedores \cite{Raynolds2007Fair/AlternativeDimensions}. A percepção dos consumidores quanto à importância do surgimento de um comércio justo contribui para o desenvolvimento sustentável ao proporcionar melhores condições de troca, gerando parcerias e a garantia dos direitos para produtores e consumidores.

O preço justo incorpora um preço mínimo aceitável ao produtor, segundo a imposição dos consumidores perante o ato de compra. Neste sentido, a percepção de preço justo credenciado pelas empresas parece depender diretamente da percepção do cliente consumidor frente à marca, por meio do reconhecimento da origem do produto e certificações obtidas. Essa premissa vincula a percepção das empresas sobre o que é realmente justo para elas ao que também e justo para os consumidores.

Para as marcas que buscam a sustentabilidade ambiental, o conceito de preço justo impacta diretamente no \textit{Brand Equity}, uma vez que, construído adequadamente, promove a afinidade do consumidor com os produtos ao compreender que está adquirindo um produto aceitável e justo pelo valor estipulado \cite{Ulian2020UnlockingAgriculture}. Pesquisas anteriores mostram que imagens planejadas relacionadas ao consumo aumentam a vontade de comprar e a satisfação de fazer negócio com a marca \cite{Schlosser2003ExperiencingIntentions}, formando conceitos positivos sobre os produtos e seu valor \cite{Chan2018NeuralTemplates}. Essa consonância obriga as empresas a criarem um posicionamento bem definido no seu nicho de atuação, permitindo assim provar sua credibilidade no que tange à atuação sustentável, podendo posteriormente cobrar um valor premium em seus produtos, trazendo margem de lucro maior em seus produtos \cite{Langen2013WhereReality}.

Quanto a gestão do posicionamento das marcas, os empreendedores afirmam utilizar diversas estratégias, tais como: (i) utilização de signos que remetem à floresta em seu design, logo, sites e embalagens; (ii) A nomenclatura dos produtos e nome distintivo utilizando, expressões idiomáticas regionais ou palavras de dialetos locais; (iii) Estratégias de \textit{Copywriting} utilizando a expressão da relação com os povos originários da Amazônia e (iv) a estratégia de diferenciação da marca focada nos valores de sustentabilidade e na causa ou propósito da conversação da floresta. A exemplo disso, podemos citar a estratégia de \textit{Product Naming}, utilizada por uma das marcas estudadas para nomear uma linha de seus produtos.

\textit{“(...) Todos os nossos rótulos são baseados em linguagem visual dos povos originários, com consentimento deles. Damos um colorido e tal, mas no geral, a arte em você é de origem deles.”} (Empreendedor 02).

\textit{“Visamos retratar as nossas parcerias através dos produtos e embalagens. Por exemplo, tem o kit expedição da Amazônia, em que representamos um caminho na mata, onde os sabores florestam e das comunidades tradicionais apresentam-se por meio das barras de chocolate. Nesse kit, tem as barras de chocolates do Vale do Jarí, dos ribeirinhos de Mocujiba, as barras Sakaguchi do cacau vindo de um agroflorestal em Tomé-Açú e a barra Xiba produzida com o cacau colhido em Barcarena, na comunidade do ribeirinho Xiba às margens do rio Acará-PA.”}  (Empreendedor 02).

\textit{“Nos nossos roteiros, gostamos de relacionar as tradicionais locais, por exemplo, temos o roteiro De Bubuia. Ficar "De Bubuia" é o jeito amazonense de dizer que quer relaxar, ficar tranquilo, ter um dia off, fazendo o que gosta, apreciando a natureza. Temos também o roteiro Perreché Tour, feito a pé (...). Essa expressão relaciona a estar descalço andar descalço.} (Empreendedor 01).

\textit{“Sim, além de símbolos, nos utilizamos as músicas regionais e nomes de animais também. É tudo pensando em coisas da região que remetam aquilo. Por exemplo, temos um lenço chamado Coá-Piranga, uma bolsa chamada Curumim, camiseta Cunhantã Vários outros produtos.”} (Empreendedor 06).

\subsubsection*{Dimensão ecológica}

No \textit{Cluster} 3, são observados relatos sobre as atividades de restauração ecológica desenvolvidas. É possível notar que para os empreendedores tanto é importante a manutenção de regiões reservadas quanto o desenvolvimento de regiões de extrativismo sustentável. Os códigos 3A, 3B e 3C, demostram que a preservação ecológica pode ser alcançada na região a partir da recuperação de locais que já foram disponíveis para atividades exploratórias, mantendo os pontos de floresta com o mínimo de interferência.

\textit{“(...) Meu pai era muito observador. Ele percebeu toda a diversidade florestal. Ele observou como era o trabalho dos polinizadores naturais e aí ele começou a ver que a floresta em pé valeria mais, assim começou a reflorestar com cacau, cupuaçu, tudo que possam imaginar de fruteiras, ou seja, de forma empírica, é o manejo sustentável, não é?”} (Empreendedor 05).

A preocupação com a dimensão ecológica se manifesta também no \textit{Cluster} 5, quando observado o código 5A, que aborda sobre de energias renováveis, a maioria dos entrevistados não dispõe de energia renovável, porém alegam ter um plano estratégico para implementação. As alegações vão desde pouco espaço, custos altos para implementação e oferta energética local a preço ainda vantajoso.

\textit{“(...) Não, ainda não, mas está na nossa agenda. Inclusive, hoje recebi a visita de uma pessoa que está fazendo um orçamento para a fábrica, para colocar, implantar um sistema, um sistema de fotovoltaica. Aí na fábrica e, inclusive, aqui em casa também.”} (Empreendedor 02).

\textit{“(...) Não, ainda não, pois ainda não tempos um espaço físico”}. (Empreendedor 01).

\textit{“(...) Não, pois utilizamos muito pouco a energia elétrica”}. (Empreendedor 04).

Quanto à redução de poluentes, os empreendimentos que recorrem a transporte, afirmaram que, sempre que possível, dão preferência a meios de transporte \textit{peer-peer}. Esses serviços permitem que as pessoas usem a mobilidade conforme a demanda necessária \cite{Shaheen2015DefinitionsUnderstanding,Santos2018SustainabilityModels}. É importante compreender que, a mobilidade compartilhada pode reduzir o congestionamento rodoviário, reduzir a demanda de infraestrutura de transporte, reduzir emissão de $CO_2$, além de custos de financiamento mais baixos em comparação com veículos proprietários \cite{Guyader2021SharedSustainability}. Quanto à redução de poluentes, os empreendimentos que usam transporte, afirmaram que, sempre que possível, dão preferência a meios de transporte \textit{peer-peer}.

\subsubsection*{Dimensão territorial}

A dimensão territorial abrange a organização do espaço e obedece a critérios superpostos de ocupação territorial e entrelaçados em uma rede natural duradoura \cite{Sachs1993EstrategiasXXI}. Tais características são definidas na concepção dos empreendedores quanto à importância de incluir áreas exploradas de paisagens sustentáveis, assim como suas espécies sustentáveis de paisagens de produção \cite{Garibaldi2004CulturalRestoration}. Eles moldam o apoio da comunidade ao projeto, o apoio da comunidade ao apoio tradicional, o apoio à ajuda para o sucesso do projeto e o apoio da comunidade \cite{Wehi2017ImportanceRestoration}. Para a visão estratégica de posicionamento de marca, essa forma de atividade preserva os locais de floresta nativa, signos atrelados ao \textit{Brand Equity}, os quais, quando somados às atitudes ambientais das empresas, influenciam diretamente a percepção do consumidor no que diz respeito à preservação ambiental desenvolvida pela empresa e auxilia no posicionamento estratégico da marca. 

Nossas percepções não apenas proporcionam aos empreendedores a oportunidade de gerar novas ideias a partir da discussão levantada sobre as dimensões da sustentabilidade, mas também uma importante discussão sobre fatores replicáveis em novos negócios de mesma natureza. 

\section{Conclusões}

O estudo quali-quantitativo permitiu demonstrar quais as percepções dos novos empreendedores da região Amazônica frente à importância da conexão da sua marca com os indicadores do desenvolvimento sustentável da floresta amazônica.

Foi possível compreender que os empreendedores percebem que a sustentabilidade depende de arranjos participativos entre a empresa, o meio ambiente e as comunidades tradicionais da região. Esta percepção se relaciona especialmente no desenvolvimento de seu \textit{Branding} focado na construção da imagem participativa e cooperativa das entidades citadas. 

Os resultados deste estudo mostram conscientização entre os empreendedors entrevistados sobre a necessidade da sustentabilidade da floresta Amazônica. Essa consciência representa um pré-requisito positivo para a implementação de estratégias de posicionamento estratégico de marcas que buscam apresentar a atividade sustentável como diferencial. 

São identificados códigos que descrevem que o atendimento aos indicadores sociais dependem da participação ativa e compartilhamento de conhecimentos ecológicos tradicionais dos povos nativos, a atuação com os agentes governamentais das reservas legais, seus bens biológicos potencialmente exploráveis sustentavelmente, a valoração dos recursos florestais pela população e uma expressiva necessidade de um maior engajamento do poder público frente ao apoio empresarial no local.

Quando observada a dimensão econômica, os comerciantes percebem dificuldades na negociação dos produtos com \textit{Brand Equity} sustentável no mercado. Este fator promove ativamente a adaptação dos empreendimentos aos mercados consumidores internacionais ambientalmente preocupados com a manutenção florestal. 

Os resultados sugerem que, para os empreendedores locais, é possível aumentar a capacidade de produção sustentável, portanto a lucratividade sem introduzir novas vinculações de recursos por meio de uma extensa valoração dos recursos sociais associados aos recursos ambientais ali existentes. Outrossim, os empreendedores afirmam que o uso de signos e certificações promovem nos seus consumidores percepções mais confiáveis sobre a atuação sustentável das marcas na floresta amazônica. 

Percebe-se que todas as partes interessadas neste desenvolvimento devem trabalhar iniciativas de promoção da sustentabilidade ambiental, a fim de criar programas socialmente justos e eficazes para a existência de longo prazo das empresas focadas no mercado, incluindo a criação de sistemas socioecológicos complexos, espaços ecológicos sustentáveis e programas mais eficazes de educação do campo.

Embora a população estudada reflita uma pequena parte de marcas atuantes na região Amazônica, identificou-se uma percepção da efetiva atuação sustentável na região, e que essa depende da conservação da floresta e seus bioinsumos, do respeito aos povos originários e seus conhecimentos tradicionais, ambos fortemente ligados a conservação da floresta.

O conhecimento direcional da marca, organizado e assertivo, promove uma melhoria na competitividade de pequenas empresas que não podem investir em Pesquisa e Desenvolvimento, demonstrando que utilizar estratégias de \textit{Branding} para aumentar a capacidade competitiva.